\subsection{Validation Methodology}
\label{sec_eval_setup}

\begin{table}[b]
 % \vspace{0.5cm}
    \centering
    \caption{Summary of evaluation sequences. 
    \#W and \#S denote the number of visible walls and distinct navigable spaces (rooms and corridors), respectively. Area, Length ($L$), and Duration ($D$) are in $m^2$, $m$, and $s$.
    CS1--CS3-Furn are distinct construction sites at different completion stages; office1 and office2 are separate buildings with multiple sequences each. }
    \begin{tabular}{l|l|c|c|c|c|c}
        \toprule
            \textbf{Env.} & \textbf{Sequence} & \textbf{\#W} & \textbf{\#S} & \textbf{Area ($\mathbf{m}^2$)} & \textbf{L (m)} & \textbf{D(s)} \\
        \midrule
        \multirow{8}{*}{\shortstack{\textbf{Offices}}} 
            & office1-1 & 10 & 4 & 100 & 90 & 670 \\
            & office1-2 & 12 & 3 & 109 & 48 & 265 \\
            & office1-3 & 12 & 4 & 113 & 99 & 524 \\
            & office1-4 & 6 & 2 & 54 & 37 & 169 \\
        \cmidrule(lr){2-7}
            & office2-1 & 10 & 3 & 168 & 68 & 268 \\
            & office2-2 & 10 & 3 & 168 & 73 & 370 \\
            & office2-3 & 10 & 3 & 70 & 41 & 172 \\
            & office2-4 & 4 & 4 & 40 & 53 & 260 \\
        \midrule
        \multirow{3}{*}{\shortstack{\textbf{Constr.}\\\textbf{Sites}}} 
            & CS1       & 16 & 5 & 48 & 51 & 377 \\
        \cmidrule(lr){2-7}
            & CS2-Mid   & 17 & 5 & 68 & 46 & 215 \\
        \cmidrule(lr){2-7}
            & CS3-Furn  & 13 & 4 & 95 & 30 & 180 \\
        \midrule
            \multicolumn{2}{c|}{\textbf{Total}} & 120 & 40 & 1033 & 642 & 3470 \\
        \bottomrule
    \end{tabular}
    \label{tbl_dataset}
\end{table}



\textbf{Datasets.}
Public indoor datasets (e.g., \cite{dataset_scannet}) typically provide RGB-D data, but to the best of our knowledge, no public dataset includes both RGB-D and the corresponding BIM, which limits the evaluation of our system. 
\add{To address this, we collected real-world sequences with the SMapper device~\cite{smapper}, a portable acquisition platform integrating an \textit{Intel RealSense D435i} RGB-D camera and an \textit{Ouster OS0-64} 3D LiDAR. ivS-Graphs uses only the RGB-D stream as input; the LiDAR point clouds are used for ground-truth trajectory generation (Section~IV-A). This type of camera is commonly mounted on robotic and handheld platforms.}

The sequences span multiple indoor layouts as summarized in Table~\ref{tbl_dataset}, including two office buildings with varying trajectories and three construction sites at different completion stages, \add{all within indoor structured environments where architectural walls are already in place}: early phase with exposed installations (\textit{CS1}), mid-phase with ongoing work (\textit{CS2-Mid}), and near-finished with furniture present (\textit{CS3-Furn}).


\begin{table}[t]
\centering
 \vspace{0.25cm}
\caption{Key Parameters and Thresholds for ivS-Graphs}
\label{tab:parameters}
\begin{tabular}{@{}lllc@{}}
\toprule
\textbf{Param.} & \textbf{Value} & \textbf{Description} & \textbf{Section} \\ \midrule
$\tau_{\perp}$ & 0.17 rad & Perpendicularity threshold for initialization & \ref{sec:wall_matching} \\
$\tau_{init}$ & 0.30 m & Max. residual for initial alignment & \ref{sec:wall_matching} \\
$\tau_p$ & 0.77 m & Plane-parameter distance threshold & \ref{sec:wall_matching} \\
$\tau_c$ & 6.00 m & Lateral centroid distance threshold & \ref{sec:wall_matching} \\
$\alpha$ & 0.70 & Weighting for plane vs. spatial proximity & \ref{sec:wall_matching} \\
$\beta$ & $10^{-2}$ & \addtwo{Base weight for wall-to-wall factors} & \ref{sec:bim_backend} \\ \bottomrule
% $\beta$ & $10^{-2}$ & Base uncertainty for wall-to-wall factors & \ref{sec:bim_backend} \\ \bottomrule
\end{tabular}
\end{table}

\textbf{Baselines.}
\change{Since no existing visual SLAM system integrates BIM priors, we evaluate against established RGB-D SLAM baselines: ORB-SLAM3~\cite{orb3}, BAD-SLAM~\cite{badslam}, and DROID-SLAM~\cite{droidslam} are selected for their reported performance and code availability.}{Since no existing visual SLAM system integrates BIM priors, we evaluate against representative RGB-D visual SLAM methods that do not rely on external architectural priors. The selected baselines cover the main RGB-D visual SLAM paradigms: feature-based (\textit{ORB-SLAM3}~\cite{orb3}), dense direct (\textit{BAD-SLAM}~\cite{badslam}), and learning-based (\textit{DROID-SLAM}~\cite{droidslam}).}
\add{Additionally, the comparison with \textit{vS-Graphs}~\cite{tourani2025vsgraphs}, one of the best-performing structure-aware visual SLAM systems that integrates wall planes into the factor graph, allows us to directly measure the effect of adding BIM constraints as our system builds on it.}


\textbf{Implementation Details.} All algorithms were executed on a workstation with an \textit{Intel Core i9-11950H} (2.60\,GHz), an \textit{NVIDIA T600 Mobile GPU}, and 32\,GB RAM. 
\add{The BIM structural data is extracted from an IFC model.
For Revit models, an automated Dynamo script converts them into IFC.
Only architectural walls are extracted; other BIM elements such as doors and windows are not used.}

\textbf{System parameters}. The parameters in Table~\ref{tab:parameters} were tuned empirically based on physical constraints: $\tau_p$ and $\tau_c$ reflect both the expected RGB-D sensor uncertainty and typical dimensions of architectural wall segments. The weighting factor $\alpha$ prioritizes geometric alignment (PPD) over spatial proximity (LCD), while $\beta$ balances BIM constraints against visual odometry, ensuring that well-matched walls exert meaningful corrections without overriding visual observations. All parameters are kept constant across every experiment and environment, demonstrating the generalization capability of the proposed configuration.

